\section{Bipartite Matching on Randomly Colored Points}

\subsection{Missing Proofs and Details of Section~\ref{sec:randomly-colored}}\label{appendix-randomly-colored}

\myparagraph{Constructing the Partial Matching.} In this part, for any cell $\cell$ of $\mcH$, we construct a matching $M'$ of $A_\cell\cup B_\cell$ that, in expectation, matches all except $\tilde{O}(n_\cell^{3/4})$ points of $B_\cell$ and has a cost $\tilde{O}(\ell_\cell n_\cell^{-1/4})$, proving Lemma~\ref{lemma:GRS-matching}.
We begin by introducing a set of notations. Let $\cell$ be any cell of $\mcH$, and let $G$ be a grid dividing $\cell$ into smaller squares.
For any square $\xi\in G$, let $A_\cell^\xi$ (resp. $B_\cell^\xi$) denote the subset of points of $A_\cell$ (resp. $B_\cell$) that lie inside $\xi$. Define the \emph{excess} of $\xi$ as $\excess(\xi):=\big||B_\cell^\xi| - |A_\cell^\xi|\big|$. Define $\excess(G):=\sum_{\xi\in\mcG}\excess(\xi)$ as the excess of $G$. We next show an important property of randomly colored point sets, which is critical in constructing the matching. 


\begin{lemma}\label{lemma:convergence_excess}
    For any square $\cell$ of $\mcH$ and a grid $G$ inside $\cell$ with cell side length $O(\ell_\cell n_\cell^{-\alpha})$, $\mbE[\excess(G)]= \tilde{O}(n_\cell^{\alpha + \frac{1}{2}})$. 
\end{lemma}
\begin{proof}
    If $\alpha\ge \frac{1}{2}$, the lemma statement holds trivially since $n_\cell^{\alpha+\frac{1}{2}} \ge n_\cell$. Therefore, we assume $\alpha\le \frac{1}{2}$.
    Using the Hoeffding’s inequality~\cite{hoeffding1994probability}, for any hypercube $\xi$ of $G$,
    \begin{equation*}
        \prob{\big||B_\cell^\xi| - |A_\cell^\xi|\big|\ge c_1\sqrt{n_\cell log n}}\le n^{-c_2}\label{eq:excess}
    \end{equation*}
    for some constants $c_1, c_2 > 1$. Since $\excess(\xi)\le n_\cell$, 
    \begin{equation*}
        \mbE[\excess(\xi)]= O(\sqrt{n_\cell\log n}).\label{eq:excess_cell}
    \end{equation*}
    Summing over all squares of $G$,
    \begin{align*}
        \mbE[\excess(G)] = \sum_{\xi\in G} \mbE[\excess(\xi)] &=O\left(\sum_{\xi\in G} \sqrt{n_\cell\log n}\right) = O\left(\sqrt{n\log n\times |G|}\right) = \tilde{O}(n^{\alpha+\frac{1}{2}}).\label{eq:excess-dense}
    \end{align*}
\end{proof}

We next use Lemma~\ref{lemma:convergence_excess} to show that there exists a low-cost high-cardinality matching inside each sub-problem.

\GRSPartialMatching*
\begin{proof}
    Let $\langle G_1, \ldots, G_t\rangle$ denote a sequence of grids, where $t=\lceil \log\log n\rceil$ and each grid $G_i$ has a side-length $O(\ell_\cell n_\cell^{-\alpha_i})$ for \[\alpha_i:=\frac{1}{2} - \frac{2^t}{2^{t+2} - 3}\left(1-\frac{1}{2^i}\right).\]
    Using these grids, we construct a matching $M'$ as follows.
    Let $A_\cell^0:=A_\cell$ and $B_\cell^0:=B_\cell$. Starting from $i=1$, we compute a matching $M_i$ from $B_\cell^{i-1}$ to $A_\cell^{i-1}$ that for each square of $G_i$, matches as many points as possible inside the square arbitrarily. Define $A_\cell^i$ (resp. $B_\cell^i$) as the set of free points of $A_\cell^{i-1}$ (resp. $B_\cell^{i-1}$) and process the next grid $G_{i+1}$. We continue this procedure until the last grid $G_t$ is processed. Define $M':=\sum_{i=1}^t M_i$. This completes the construction of $M'$. In the following, we first show that the total number of free points with respect to $M'$ is $O(n_\cell^{3/4})$; we then show that $w(M')=O(\ell_\cell^2n_\cell^{1/4})$ and conclude the lemma statement.

    The matching $M$ matches as many points as possible inside each square of the grid $G_t$. Therefore, the total number of unmatched points is equal to the excess of $G_t$, which by Lemma~\ref{lemma:convergence_excess} is 
    \begin{align}
        \mbE[\excess(G_t)] &= \tilde{O}(n_\cell^{\alpha_t+\frac{1}{2}}) = O(n_\cell^{\frac{3}{4}+\frac{1}{4(2^{t+2}-3)}}) = O(n_\cell^{\frac{3}{4}+\frac{1}{16\log n-12}}) = O(n_\cell^{\frac{3}{4}}).
    \end{align}
    We next analyze the expected cost of $M'$. By the linearity of expectation,
    \begin{equation}
        \mbE[w(M')] = \sum_{i=1}^t \mbE[w(M_i)].
        \label{eq:expected_cost_1}
    \end{equation}
    For $i=1$, since all matching edges in $M_1$ have a squared Euclidean cost at most $O((\ell_\cell n_\cell^{-\alpha_1})^2)$, the cost of $M_1$ would be 
    \begin{equation}
        w(M_1) = O(n_\cell\times (\ell_\cell n_\cell^{-\alpha_1})^2) = O(\ell_\cell^2 n_\cell^{\frac{2^t}{2^{t+2}-3}}) = O(\ell_\cell^2 n_\cell^{\frac{1}{4}}).\label{eq:expected_cost_2}
    \end{equation}
    Finally, for each $1<i\le t$, the matching $M_i$ consists of matching edges with squared Euclidean cost of at most $O((\ell_\cell n_\cell^{-\alpha_i})^2)$. By Lemma~\ref{lemma:convergence_excess}, the expected number of matching edges in $M_i$ is at most $\tilde{O}(n_\cell^{\alpha_{i-1}+\frac{1}{2}})$; therefore,
    \begin{equation}
        \mbE[w(M_i)] = \tilde{O}(n_\cell^{\alpha_{i-1}+\frac{1}{2}}\times (\ell_\cell n_\cell^{-\alpha_i})^2) = \tilde{O}(\ell_\cell^2 n_\cell^{\frac{2^t}{2^{t+2}-3}}) = \tilde{O}(\ell_\cell^2 n_\cell^{\frac{1}{4}}).\label{eq:expected_cost_3}
    \end{equation}
    Combining Equations~\eqref{eq:expected_cost_1},~\eqref{eq:expected_cost_2}, and~\eqref{eq:expected_cost_3}, $\mbE[w(M)] = \tilde{O}(\ell_\cell^2 n_\cell^{\frac{1}{4}})$.
\end{proof}

\paragraph{Bounding the number of iterations.} We next show that the number of iterations of the merge step and the extended Hungarian search procedure on any cell $\cell$ of $\mcH$ is bounded by $O(n^{3/4})$.


\matchingIterations*
\begin{proof}
    From Lemma~\ref{lemma:local-inactive}, the total number of iterations of the merge step on $\cell$ is $O(|B_\cell^\mcC|)$.
    Next, note that the parameter $k$ is set to $0$; therefore, there are no free points inside $\cell$ when the merge step is executed on $\cell$. When erasing the divider of $\cell$, our algorithm creates $|B_\cell^\mcC|$ free points. Each iteration of the merge step either (i) finds an augmenting path, which reduces the number of free points by one, or (ii) finds an alternating path to a matched point $b\in B_{\cell}$ with $y(b)=\ymax$, which does not change the number of free points inside $\cell$. Therefore, after the merge step on $\cell$, the number of free points remaining inside $\cell$ is at most $|B_\cell^\mcC|$, and our algorithm executes one extended Hungarian search procedure for each remaining free point.
\end{proof}

\GRSIters*
\begin{proof}
Let $M'$ denote the matching constructed in Lemma~\ref{lemma:GRS-matching}, and let $M$ denote the matching of the extended matching $M^\mcC$ maintained by our algorithm. Let $M_\cell$ denote the matching edges of $M$ that lie inside $\cell$. Note that by lemma~\ref{lemma:no-cross}, no matching edges can cross the boundaries of $\mcC$.
Let $\mcP_{\mathrm{aug}}$ (resp. $\mcP_{\mathrm{alt}}$) denote the set of augmenting paths (resp. alternating paths) with an endpoint in $B_\cell^\mcC$ in the symmetric difference $M_\cell\oplus M'$. Clearly, $|B_\cell^\mcC|=|\mcP_{\mathrm{alt}}| + |\mcP_{\mathrm{aug}}|$. For the alternating paths, $|\mcP_{\mathrm{alt}}|=O(n^{3/4})$ since each alternating path has one free endpoint in $M'$. 
Next, we show that $|\mcP_{\mathrm{aug}}|=O(n^{3/4})$.

As shown in Equation~\eqref{eq:net-cost-sum-2},
\begin{align}
    \sum_{b\in B_{\mathrm{aug}}} y(b) = \sum_{P\in \mcP_{\mathrm{aug}}} y(b_P)\le \sum_{P\in \mcP_{\mathrm{aug}}} \phi(P) \le w(M').\label{eq:net-cost-sum-1-m}
\end{align}

Define $\alpha:=\ell_\cell n^{-1/4}$. Each free point $b\in B_{\mathrm{aug}}$ is called a \emph{close} (resp. \emph{far}) point if the Euclidean distance of $b$ to the divider of $\cell$ is at most (resp. more than) $\alpha$. Let $B_{\mathrm{aug}}^{\mathrm{close}}$ (resp. $B_{\mathrm{aug}}^{\mathrm{far}}$) denote the set of all close (resp. far) points of $B_{\mathrm{aug}}$. By Lemma~\ref{lemma:margin},
\begin{equation}\label{eq:close-points-m}
    |B_{\mathrm{aug}}^{\mathrm{close}}| = O(n_\cell n^{-1/4}) = O(n^{3/4}).
\end{equation}
For each far point $b\in B_{\mathrm{aug}}^{\mathrm{far}}$, $y(b)=\distance{b}{\divider\cell}\ge \alpha^2$ since $b$ is matched to the divider $\divider\cell$ (note that the distance function $\distance{\cdot}{\cdot}$ is the squared Euclidean distance). Therefore, 
\begin{equation}\label{eq:net-cost-3-m}
    \sum_{b\in B_{\mathrm{aug}}} y(b) \ge \sum_{b\in B_{\mathrm{aug}}^{\mathrm{far}}} y(b) \ge \alpha^2\times |B_{\mathrm{aug}}^{\mathrm{far}}|.
\end{equation}
Combining Equations~\eqref{eq:net-cost-sum-1-m} and~\eqref{eq:net-cost-3-m},
\begin{equation}\label{eq:far-points-m}
    |B_{\mathrm{aug}}^{\mathrm{far}}| \le \frac{\sum_{b\in B_{\mathrm{aug}}} y(b)}{\alpha^2} \le \frac{w(M')}{\alpha^2} = O(n^{3/4}).
\end{equation}
Combining with Equation~\eqref{eq:close-points-m},
\begin{align*}
    |B_\cell^\mcC| &\le |B_{\mathrm{aug}}| + |B_{\mathrm{alt}}| \le |B_{\mathrm{aug}}^{\mathrm{close}}| + |B_{\mathrm{aug}}^{\mathrm{far}}| + |B_{\mathrm{alt}}| = \tilde{O}(n^{3/4}).
\end{align*}
\end{proof}





\subsection{Analysis for General \texorpdfstring{$d$}{} and \texorpdfstring{$q$}{}}\label{sec:appendix-GRS}
For a point set $U$ of $2n$ points in $d$ dimensions and any $q\ge 1$, let $A$ denote a random subset of $n$ points of $U$ and let $B=U\setminus A$. 
To compute a minimum-cost perfect matching between $A$ and $B$ under $\ell_2^q$ distances, we construct our hierarchical partitioning $\mcH$ with a parameter $\lambda=9n^{-\frac{1}{d+2}}$ and execute our algorithm from Section~\ref{sec:k-seq} by setting $k=0$. In the following, we extend our analysis from Section~\ref{sec:randomly-colored} to any dimension $d\ge 2$ and any $q\ge 1$ and show the following result.


\begin{restatable}{lemma}{GRSgeneral}\label{lemma:GRS-matching-general}
    Suppose $U$ is a set of $2n$ points inside the unit $d$-dimensional hypercube, $d\ge 2$, and $A$ is a subset chosen uniformly at random from all subsets of size $n$. Let $B=U\setminus A$. Then, for any parameters $q\ge 1$, the expected running time of our algorithm for computing the minimum-cost matching on the complete bipartite graph on $A$ and $B$ under $\ell_2^q$ costs is
    \begin{equation*}
    \begin{cases}
        \tilde{O}(n^{2-\frac{q}{(q+1)d}}\Phi(n)\log \Delta),\qquad &q\le\frac{d}{2},\\
        \tilde{O}(n^{2-\frac{1}{d+2}}\Phi(n)\log \Delta),\qquad &q>\frac{d}{2}.
    \end{cases}
\end{equation*}
\end{restatable}

% As in Section~\ref{sec:improved-analysis}, we begin by constructing a partial matching that has a high cardinality and matches
Note that when $q\le \frac{d}{2}$, the bound claimed in Lemma~\ref{lemma:GRS-matching-general} is identical to the runtime bound proven in~\cite[Theorem B.2]{gattani2023robust} and their analysis directly applies to our algorithm; hence, in the following, we analyze the running time of our algorithm assuming $q>\frac{d}{2}$.

We begin by showing that there exists a low-cost high-cardinality partial matching.

\begin{lemma}\label{lemma:GRS-matching-d}
    For any cell $\cell$ of $\mcH$, there exists a matching $M'$ that, in expectation, matches all except $O(n_\cell^{1-\frac{1}{d+2}})$ points of $B_\cell$ and has a cost $O((2\ell_\cell)^q n_\cell^{1-\frac{q+1}{d+2}})$.
\end{lemma}
\begin{proof}
Define $r:= \frac{2q}{d}>1$ and $t:=\lceil \log_r \log_2 n\rceil$. Let $\beta:=q(d+2)r^t - d(q+1)$. Define a sequence of grids $\langle G_1, \ldots, G_t\rangle$, where each grid $G_i$ has a side-length $O(\ell_\cell n^{-\alpha_i})$ for \[\alpha_i:=\frac{1}{d} - \frac{r^{t+1}}{\beta}\left(1-\frac{1}{r^i}\right).\]

We construct a matching $M$ as follows.
Let $A_\cell^0:=A_\cell$ and $B_\cell^0:=B_\cell$. Starting from $i=1$, let $M_i$ be a matching from $B_\cell^{i-1}$ to $A_\cell^{i-1}$ that matches as many points as possible inside each hypercube of $G_i$. Let $A_\cell^i$ and $B_\cell^i$ denote the set of free points of $A_\cell^{i-1}$ and $B_\cell^{i-1}$, respectively. Set $i\leftarrow i+1$ and continue this procedure until the last grid $G_t$ is processed. Define $M:=\sum_{i=1}^t M_i$. 
In the following, we show that the total number of free points with respect to $M$ is $\tilde{O}(n_\cell^{1-\frac{1}{d+2}})$ and the cost of $M$ is $\mbE[w(M)]=O((2\ell_\cell)^qn_\cell^{1-\frac{q+1}{d+2}})$.

The matching $M$ matches as many points as possible inside each square of the grid $G_t$. Therefore, the total number of unmatched points is equal to the excess of $G_t$, which by Lemma~\ref{lemma:convergence_excess-d} below is 
\begin{align}
    \mbE[\excess(G_t)] = \tilde{O}(n_\cell^{\frac{d}{2}\alpha_t+\frac{1}{2}}) = \tilde{O}(n_\cell^{1 - \frac{1}{d+2}}), \label{eq:expected_free_general}
\end{align}
where the second equality holds since
\begin{align*}
    \frac{d}{2}\alpha_t+\frac{1}{2} &= \frac{d}{2}\left(\frac{1}{d} - \frac{r^{t+1}}{\beta}(1-\frac{1}{r^t}) \right) + \frac{1}{2}= \left(\frac{1}{2} - \frac{qr^{t}-q}{q(d+2)r^t - d(q+1)} \right) + \frac{1}{2}\\ &= 1 - \frac{1}{d+2} + \frac{q-\frac{d(q+1)}{d+2}}{q(d+2)r^t - d(q+1)} \le 1 - \frac{1}{d+2} + \frac{1}{2\log n}. 
\end{align*}
We next analyze the expected cost of $M$. For $i=1$, since all matching edges in $M_1$ has a cost of at most $O((\ell_\cell n_\cell^{-\alpha_1})^q)$, the cost of $M_1$ would be 
\begin{equation}
    w(M_1) = O(n_\cell\times (\ell_\cell n_\cell^{-\alpha_1})^q) = O((2\ell_\cell)^q n_\cell^{1-\frac{q+1}{d+2}}),\label{eq:expected_cost_2_general}
\end{equation}
where the second equality holds since
\begin{align*}
    1-q\alpha_1 &= 1-\frac{q}{d}+\frac{q(r-1)r^t}{\beta}= 1-\frac{q}{d}+\frac{\frac{2q^2-qd}{d}r^t}{q(d+2)r^t - d(q+1)} \\ &= 1-\frac{q}{d}+\frac{2q-d}{d(d+2)} + \frac{\frac{(2q-d)(q+1)}{d+2}}{q(d+2)r^t - d(q+1)}\le 1-\frac{q+1}{d+2} + \frac{q}{d\log n}.
\end{align*}
Finally, for each $1<i\le t$, the matching $M_i$ consists of matching edges with a cost of $O((\ell_\cell n_\cell^{-\alpha_i})^q)$. By Lemma~\ref{lemma:convergence_excess-d}, the expected number of matching edges in $M_i$ is at most $\tilde{O}(n_\cell^{\frac{d}{2}\alpha_{i-1}+\frac{1}{2}})$; therefore,
\begin{equation}
    \mbE[w(M_i)] = \tilde{O}(n_\cell^{\frac{d}{2}\alpha_{i-1}+\frac{1}{2}}\times (\ell_\cell n_\cell^{-\alpha_i})^q) = \tilde{O}((2\ell_\cell)^q n_\cell^{1-\frac{q+1}{d+2}}),\label{eq:expected_cost_3_general}
\end{equation}
where the second equality is resulted from as follows.
\begin{align*}
    \frac{d}{2}\alpha_{i-1} + \frac{1}{2}-q\alpha_i &= \left(\frac{1}{2}-\frac{qr^t(1-\frac{1}{r^{i-1}})}{\beta}\right) + \frac{1}{2} - \left(\frac{q}{d}-\frac{qr^{t+1}(1-\frac{1}{r^{i}})}{\beta}\right)\\ &=  1-\frac{q}{d}+\frac{q(r-1)r^t}{\beta}\le 1-\frac{q+1}{d+2} + \frac{q}{d\log n}.
\end{align*}
Combining Equations~\eqref{eq:expected_cost_2_general} and~\eqref{eq:expected_cost_3_general}, the expected cost of $M$ is \[\mbE[w(M)]=\sum_{i=1}^t\mbE[w(M_i)]=O((2\ell_\cell)^qn_\cell^{1-\frac{q+1}{d+2}}).\]
\end{proof}

\begin{lemma}\label{lemma:convergence_excess-d}
    For any random partitioning of a set $U$ of $2n$ points inside the $d$-dimensional unit hypercube into two sets $A$ and $B$ of $n$ points each, a hypercube $\cell$ of $\mcH$, and a grid $G$ inside $\cell$ with cell side length $O(\ell_\cell n_\cell^{-\alpha})$, $\mbE[\excess(G)]= \tilde{O}(n_\cell^{\frac{d}{2}\alpha + \frac{1}{2}})$. 
\end{lemma}
\begin{proof}
    For $\alpha\ge \frac{1}{d}$, the lemma statement holds trivially since $n_\cell^{\frac{d}{2}\alpha+\frac{1}{2}} \ge n_\cell$. Therefore, we assume $\alpha\le \frac{1}{d}$.
    Using the Hoeffding’s inequality~\cite{hoeffding1994probability}, for any hypercube $\xi$ of $G$,
    \begin{equation*}
        \prob{|B_\cell^\xi| - |A_\cell^\xi|\ge c_1\sqrt{n_\cell log n}}\le n^{-c_2}\label{eq:excess-d}
    \end{equation*}
    for some constants $c_1, c_2 > 1$. Since $\excess(\xi)\le n_\cell$, 
    \begin{equation*}
        \mbE[\excess(\xi)]= O(\sqrt{n_\cell\log n}).
    \end{equation*}
    Summing over all squares of $G$,
    \begin{align*}
        \mcB[\excess(G)] = \sum_{\xi\in G} \mbE[\excess(\xi)] &=O\left(\sum_{\xi\in G} \sqrt{n_\cell\log n}\right) = O\left(\sqrt{n\log n\times |G|}\right) = \tilde{O}(n^{\frac{d}{2}\alpha+\frac{1}{2}}).\label{eq:excess-dense-d}
    \end{align*}
\end{proof}


We use Lemma~\ref{lemma:GRS-matching-d} to show that the number of points in $B_\cell^\mcC$ is $\tilde{O}(n^{1-\frac{1}{d+2}})$. Let $M'$ denote the matching constructed in Lemma~\ref{lemma:GRS-matching-d}, and let $M$ denote the matching of the extended matching $M^\mcC$ maintained by our algorithm. Let $M_\cell$ denote the matching edges of $M$ that lie inside $\cell$. 
Let $\mcP_{\mathrm{aug}}$ (resp. $\mcP_{\mathrm{alt}}$) denote the set of augmenting paths (resp. alternating paths) with an endpoint in $|B_\cell^\mcC|$ in the symmetric difference $M_\cell\oplus M'$.
As discussed above, $|\mcP_{\mathrm{alt}}|\le |F(M')|=O(n^{1-\frac{1}{d+2}})$. 
We partition the free endpoints of $\mcP_{\mathrm{aug}}$ into the set $B^{\mathrm{close}}_{\mathrm{aug}}$ that are at a Euclidean distance closer than $\lambda'_\cell=2\ell_\cell n^{-\frac{1}{d+2}}$ to the divider $\divider\cell$ of $\cell$ and the set $B^{\mathrm{far}}_{\mathrm{aug}}$ that are at a Euclidean distance further than $\lambda'_\cell$ from $\divider\cell$. From Lemma~\ref{lemma:margin},  $|B^{\mathrm{close}}_{\mathrm{aug}}|=\tilde{O}(n^{1-\frac{1}{d+2}})$. Finally, by Equation~\eqref{eq:net-cost-sum-1-m}, $\sum_{b\in B^{\mathrm{far}}_{\mathrm{aug}}}y(b) \le w(M')$. Since the dual weight of each free point in $B^{\mathrm{far}}_{\mathrm{aug}}$ is at least $(\lambda'_\cell)^q$, we get a bound of $\tilde{O}(n^{1-\frac{1}{d+2}})$ on the number of such points. 

Therefore, the total number of iterations of the merge and extended Hungarian search processes on each cell $\cell$ is $\tilde{O}(n^{1-\frac{1}{d+2}})$. Each iteration of the merge step and the search procedure takes $\tilde{O}(n_\cell\Phi(n))$ time; as a result, the total execution time of our algorithm on $\cell$ would be $\tilde{O}(n^{1-\frac{1}{d+2}}n_\cell\Phi(n))$. Since each point participates in $O(\log(n\Delta))$ cells in $\mcH$, the total execution time of our algorithm across all cells would be $\tilde{O}(n^{2-\frac{1}{d+2}}\Phi(n)\log\Delta)$, proving Lemma~\ref{lemma:GRS-matching-general}.

% The rest of the proof 

% Let $M_\cell$ denote the matching computed at the beginning of the conquer step of the GRS algorithm when processing $\cell$.
% Similar to Section~\ref{sec:improved-analysis}, we bound $|F(M_\cell)|$ as $|F(M_\cell)\setminus F(M)| + |F(M)|$, where by Equation~\eqref{eq:expected_free_general}, $\mbE[|F(M)|]=\tilde{O}(n^{1-\frac{1}{d+2}})$.  In the following, we show that $\mbE[|F(M_\cell)\setminus F(M)|]$ is also $\tilde{O}(n^{1-\frac{1}{d+2}})$ and conclude that the conquer step of the GRS algorithm when processing $\cell$ executes $\tilde{O}(n^{1-\frac{1}{d+2}})$ iterations. Note that each iteration takes $\tilde{O}(n_\cell\Phi(n_\cell))$ time, where for each level $i$, $\sum_{\cell'\in\mcL[i]}n_\cell\le n$; recall that $\mcL[i]$ denotes the set of all hypercubes of the quadtree at level $i$. Therefore, processing all hypercubes of each level $i$ of quadtree takes $\tilde{O}(n^{2-\frac{1}{d+2}}\Phi(n))$ time. Summing over all $O(\log \Delta)$ levels of the quadtree, we obtain the bounds claimed in Lemma~\ref{lemma:GRS-matching-general}. 


% To bound $\mbE[|F(M_\cell)\setminus F(M)|]$, similar to Section~\ref{sec:improved-analysis} part (iii), we categorize the free points in $F(M_\cell)$ as close and far as follows. Define $\theta:=2\ell_\cell n^{-\frac{1}{d+2}}$. 
% For each point $b\in B_\cell$, recall that $\cell_b\in\children\cell$ denotes the child of $\cell$ containing $b$. Each free point $b\in F(M_\cell)$ is called a \emph{close} (resp. \emph{far}) point if the Euclidean distance of $b$ to the boundaries of $\cell_b$ is at most (resp. more than) $\theta$, and let $X$ (resp. $Y$) denote the set of all close (resp. far) points of $F(M_\cell)\setminus F(M)$. By~\cite[Lemma~2.1]{gattani2023robust}, the expected number of close points of $B_\cell$ is 
% \begin{equation}\label{eq:close-points_general}
%     \mbE[|X|] = O(n^{-\frac{1}{d+2}}\mbE[n_\cell]) = O(n^{1-\frac{1}{d+2}}).
% \end{equation}
% For each far point $b\in F(M_\cell)\setminus F(M)$, since $M_\cell, y(\cdot)$ obtained from combining $\cell'$-optimal matching for all children $\cell'\in\children\cell$, $y(b)\ge \theta^q$; therefore, 
% \begin{equation}\label{eq:net-cost-3_general}
%     \sum_{b\in F(M_\cell)\setminus F(M)} y(b) \ge \sum_{b\in Y} y(b) \ge \theta^q\times |Y|.
% \end{equation}
% Combining Equations~\eqref{eq:net-cost-sum-2} and~\eqref{eq:net-cost-3_general},
% \begin{equation}\label{eq:far-points_general}
%     \mbE[|Y|] \le \frac{\mbE[\sum_{b\in F(M_\cell)\setminus F(M)} y(b)]}{\theta^q} \le \frac{\mbE[w(M)]}{\theta^q} = O(n^{1-\frac{1}{d+2}}).
% \end{equation}
% Combining with Equation~\eqref{eq:close-points_general}, $\mbE[|F(M_\cell)\setminus F(M)|]=\tilde{O}(n^{1-\frac{1}{d+2}})$, as desired.

% Note that Lemma~\ref{lemma:GRS-analysis-p2} follows by simply setting $d=2$ in Lemma~\ref{lemma:GRS-matching-general}. 